\chapter{Quantum Singleton Bound}
\label{chap:quantum-singleton}
%%%%%%%%%%%%%%%%%%%%%%%%%%%%%%%%%%%%%%%%%%%%%%%%%%%%%%%%%%%%%%%%%%%%%%%%%%%%%%%
\section{Overview}

This chapter proves the \textbf{quantum Singleton bound}:
for a stabilizer code with logical dimension $k$ and distance $d$
on $n$ qudits over a prime field $\mathbb{F}_p$,
\[
k + 2(d-1) \;\le\; n.
\]
The proof uses the symplectic vector space $V = \mathbb{F}_p^n \times \mathbb{F}_p^n$
with the standard symplectic form, and the key idea of erasure correctability.

%%%%%%%%%%%%%%%%%%%%%%%%%%%%%%%%%%%%%%%%%%%%%%%%%%%%%%%%%%%%%%%%%%%%%%%%%%%%%%%
\section{Symplectic vector space}

\begin{definition}[Symplectic form]
\lean{sym_form}
\label{sym_form}
\leanok
\statementsource{quant-ph.9705052}{
  pdf-e658890bd7ef-p007-b004,
  pdf-e658890bd7ef-p009-b007
}
\statementsource{arXiv.2602.20186}{pdf-585ff62e2309-p002-b007}
\statementsource{quant-ph.9608006}{
  pdf-4fa1f624d1c5-p002-b005,
  pdf-4fa1f624d1c5-p003-b003
}
\statementsource{quant-ph.0508070}{pdf-4bcdc0ee9c86-p005-b004}
\statementsource{quant-ph.0005008}{
  pdf-1ac7bb3b99cc-p003-b003,
  pdf-1ac7bb3b99cc-p004-b002
}
For $u = (x,z), v = (x',z') \in V = \mathbb{F}_p^n \times \mathbb{F}_p^n$,
\[
\omega(u,v) \;=\; \sum_{i=0}^{n-1}(x_i z'_i - z_i x'_i).
\]
\end{definition}

\begin{lemma}[Additivity of the symplectic form in the first argument]
\lean{sym_form_add_left}
\label{sym_form_add_left}
\statementsource{arXiv.2602.20186}{pdf-585ff62e2309-p002-b007}
\lean{sym_form_add_right}
\label{sym_form_add_right}
\lean{sym_form_smul_left}
\label{sym_form_smul_left}
\lean{sym_form_smul_right}
\label{sym_form_smul_right}
\leanok
\uses{sym_form}
\difficulty{3}
Let $p$ be prime and work over the space $V = \bbf_p^n \times \bbf_p^n$ equipped with
the symplectic form $\omega(u,v) = \sum_{i=0}^{n-1}(x_i z'_i - z_i x'_i)$, where $u =
(x,z)$ and $v = (x',z')$. Then for all $x, y, z \in V$,
\[
\omega(x + y, z) = \omega(x, z) + \omega(y, z).
\]

Fix a prime $p$ and a natural number $n$, and let $V = \bbf_p^n \times \bbf_p^n$, where
$\bbf_p = \bbz/p\bbz$. Equip $V$ with the symplectic form $\omega(u,v) =
\sum_{i=0}^{n-1}(x_i z'_i - z_i x'_i)$ for $u = (x,z)$ and $v = (x',z')$. Then $\omega$
is additive in its second argument: for all $x, y, z \in V$, \[ \omega(x,\, y + z) =
\omega(x,y) + \omega(x,z). \]

Let $p$ be prime and work over the prime field $\bbf_p$, with $V = \bbf_p^n \times
\bbf_p^n$ carrying the symplectic form $\omega$ given by $\omega(u,v) =
\sum_{i=0}^{n-1}(x_i z'_i - z_i x'_i)$ for $u=(x,z)$ and $v=(x',z')$. Then for every
scalar $c \in \bbf_p$ and all $x, y \in V$,
\[
\omega(c\,x,\, y) \;=\; c\,\omega(x, y),
\]
where $c\,x$ denotes the coordinatewise scalar multiplication of $x$ by $c$.

Let $p$ be a prime, and equip $V = \bbf_p^n \times \bbf_p^n$ with the symplectic form
$\omega(u,v) = \sum_{i=0}^{n-1}(x_i z'_i - z_i x'_i)$, where $u = (x,z)$ and $v =
(x',z')$. Then for every scalar $c \in \bbf_p$ and all $x, y \in V$,
\[
\omega(x, c\,y) \;=\; c\,\omega(x, y),
\]
where $c\,y$ denotes componentwise scalar multiplication.
\end{lemma}

\begin{lemma}[Antisymmetry of the symplectic form]
\lean{sym_form_swap}
\label{sym_form_swap}
\leanok
\proofsource{arXiv.2602.20186}{pdf-585ff62e2309-p002-b007}
\statementsource{arXiv.2602.20186}{pdf-585ff62e2309-p002-b007}
\uses{sym_form}
\difficulty{2}
Fix a prime $p$ and an integer $n$, and equip $V = \bbf_p^n \times \bbf_p^n$ with the
bilinear form $\omega(u,v) = \sum_{i=0}^{n-1}(x_i z'_i - z_i x'_i)$, where $u = (x,z)$
and $v = (x',z')$. Then $\omega$ is antisymmetric: for all $u, v \in V$, \[ \omega(u,v)
= -\,\omega(v,u). \]
\end{lemma}

\begin{lemma}[Nondegeneracy of the symplectic form]
\lean{sym_form_nondegenerate}
\label{sym_form_nondegenerate}
\leanok
\proofsource{arXiv.2602.20186}{
  pdf-585ff62e2309-p003-b001,
  pdf-585ff62e2309-p003-b002
}
\uses{sym_form}
\difficulty{4}
Let $p$ be a prime and $n \ge 0$, and let $\omega$ be the symplectic form on $V =
\bbf_p^n \times \bbf_p^n$. If a vector $u \in V$ satisfies $\omega(u,v) = 0$ for every
$v \in V$, then $u = 0$.
\end{lemma}

\begin{definition}[Bundled bilinear form]
\lean{symB}
\label{symB}
\leanok
\statementsource{arXiv.2602.20186}{pdf-585ff62e2309-p002-b007}
\uses{sym_form}
$\omega$ packaged as a \texttt{LinearMap.BilinForm}.
\end{definition}

%%%%%%%%%%%%%%%%%%%%%%%%%%%%%%%%%%%%%%%%%%%%%%%%%%%%%%%%%%%%%%%%%%%%%%%%%%%%%%%
\section{Support, weight, and isotropic subspaces}

\begin{definition}[Support and weight]
\lean{supp}
\label{supp}
\statementsource{quant-ph.9705052}{pdf-e658890bd7ef-p002-b002}
\statementsource{arXiv.2602.20186}{pdf-585ff62e2309-p003-b005}
\statementsource{quant-ph.9608006}{pdf-4fa1f624d1c5-p002-b005}
\statementsource{quant-ph.0508070}{pdf-4bcdc0ee9c86-p003-b007}
\statementsource{quant-ph.0005008}{pdf-1ac7bb3b99cc-p004-b003}
\lean{wt}
\label{wt}
\leanok
The \emph{support} $\mathrm{supp}(v) \subseteq \mathrm{Fin}\,n$ consists of
coordinates where either component of $v$ is nonzero;
the \emph{weight} is $\mathrm{wt}(v) = |\mathrm{supp}(v)|$.
\end{definition}

\begin{lemma}[Weight bound]
\lean{wt_le_n}
\label{wt_le_n}
\leanok
\uses{wt}
\difficulty{2}
Fix a prime $p$ and a natural number $n$, and let $v \in \bbf_p^n \times \bbf_p^n$. Then
the weight $\mathrm{wt}(v)$ satisfies $\mathrm{wt}(v) \le n$.
\end{lemma}

\begin{definition}[Support submodule]
\lean{V_sub}
\label{V_sub}
\leanok
\statementsource{arXiv.2602.20186}{pdf-585ff62e2309-p003-b006}
For $C \subseteq \mathrm{Fin}\,n$, the \emph{support submodule}
$V_C = \{v \in V \mid \mathrm{supp}(v) \subseteq C\}$.
\end{definition}

\begin{lemma}[Dimension of the support submodule]
\lean{dim_V_sub}
\label{dim_V_sub}
\leanok
\proofsource{arXiv.2602.20186}{
  pdf-585ff62e2309-p003-b007,
  pdf-585ff62e2309-p003-b008
}
\uses{V_sub}
\difficulty{1}
Fix a prime $p$ and an integer $n$, and let $V = \bbf_p^{\,n} \times \bbf_p^{\,n}$ be
the space of pairs of coordinate vectors over the prime field $\bbf_p$. For a subset $C
\subseteq \{1, \dots, n\}$, let $V_C \subseteq V$ be the support submodule consisting of
those pairs $(u, w)$ for which $u_i = 0$ and $w_i = 0$ at every index $i \notin C$. Then
$V_C$ has dimension $\dim_{\bbf_p}(V_C) = 2\lvert C\rvert$ over $\bbf_p$.
\end{lemma}

\begin{definition}[Symplectic orthogonal complement]
\lean{sym_orth}
\label{sym_orth}
\leanok
\statementsource{quant-ph.9705052}{pdf-e658890bd7ef-p005-b003}
\statementsource{quant-ph.9608006}{pdf-4fa1f624d1c5-p003-b004}
\statementsource{quant-ph.0508070}{pdf-4bcdc0ee9c86-p005-b004}
\statementsource{quant-ph.0005008}{pdf-1ac7bb3b99cc-p002-b001}
\uses{sym_form}
For a submodule $S \le V$,
$S^{\perp_\omega} = \{v \in V \mid \forall s \in S,\; \omega(v,s) = 0\}$.
\end{definition}

\begin{lemma}[Dimension of the symplectic orthogonal complement]
\lean{finrank_sym_orth}
\label{finrank_sym_orth}
\leanok
\uses{sym_orth, sym_form_nondegenerate}
\difficulty{2}
Let $p$ be a prime and equip the $2n$-dimensional space $V = \bbf_p^{\,n} \times
\bbf_p^{\,n}$ over the prime field $\bbf_p$ with the symplectic form
$\omega\big((x,z),(x',z')\big) = \sum_{i=1}^{n}(x_i z'_i - z_i x'_i)$. Then for every
subspace $S \le V$, its symplectic orthogonal complement $S^{\perp_\omega} = \{v \in V
\mid \omega(v,s) = 0 \text{ for all } s \in S\}$ has dimension \[ \dim_{\bbf_p}
S^{\perp_\omega} = 2n - \dim_{\bbf_p} S. \]
\end{lemma}

\begin{definition}[Isotropic submodule]
\lean{IsIsotropic}
\label{IsIsotropic}
\leanok
\statementsource{quant-ph.9705052}{
  pdf-e658890bd7ef-p004-b008,
  pdf-e658890bd7ef-p007-b004
}
\statementsource{arXiv.2602.20186}{pdf-585ff62e2309-p003-b010}
\statementsource{quant-ph.9608006}{pdf-4fa1f624d1c5-p003-b004}
\statementsource{quant-ph.0508070}{pdf-4bcdc0ee9c86-p005-b005}
\statementsource{quant-ph.0005008}{
  pdf-1ac7bb3b99cc-p001-b007,
  pdf-1ac7bb3b99cc-p004-b002
}
\uses{sym_form}
$S$ is \emph{isotropic} if $S \le S^{\perp_\omega}$, i.e.\ $\omega(u,v) = 0$
for all $u,v \in S$.
\end{definition}

%%%%%%%%%%%%%%%%%%%%%%%%%%%%%%%%%%%%%%%%%%%%%%%%%%%%%%%%%%%%%%%%%%%%%%%%%%%%%%%
\section{Code parameters and erasure correctability}

\begin{definition}[Quantum code distance]
\lean{code_dist}
\label{code_dist}
\leanok
\statementsource{quant-ph.9705052}{
  pdf-e658890bd7ef-p005-b005,
  pdf-e658890bd7ef-p002-b006
}
\statementsource{arXiv.2602.20186}{pdf-585ff62e2309-p003-b010}
\statementsource{quant-ph.9608006}{
  pdf-4fa1f624d1c5-p004-b001,
  pdf-4fa1f624d1c5-p002-b006
}
\statementsource{quant-ph.0508070}{pdf-4bcdc0ee9c86-p005-b005}
\statementsource{quant-ph.0005008}{
  pdf-1ac7bb3b99cc-p004-b003,
  pdf-1ac7bb3b99cc-p005-b002
}
\statementsource{PhysRevA.55.900}{pdf-bbdd8e5c3949-p008-b003}
\uses{wt, sym_orth, IsIsotropic}
$d(S) = \min\{\mathrm{wt}(v) \mid v \in S^{\perp_\omega} \setminus S,\; \mathrm{wt}(v) \neq 0\}$
(or $0$ if no such $v$ exists).
\end{definition}

\begin{lemma}[Code distance is at most the block length]
\lean{code_dist_le_n}
\label{code_dist_le_n}
\leanok
\uses{code_dist, wt_le_n}
\difficulty{4}
Fix a prime $p$ and a natural number $n$, and let $V = \bbf_p^{\,n} \times \bbf_p^{\,n}$
carry the symplectic form $\omega\bigl((x,z),(x',z')\bigr) = \sum_{i} (x_i z'_i - z_i
x'_i)$. For every $\bbf_p$-subspace $S \le V$, the associated code distance $d(S)$ — the
least weight $\mathrm{wt}(v) = |\mathrm{supp}(v)|$ attained by a vector $v$ lying in the
symplectic orthogonal complement $S^{\perp_\omega}$ but not in $S$ — satisfies $d(S) \le
n$.
\end{lemma}

\begin{definition}[Erasure correctability]
\lean{correctable}
\label{correctable}
\leanok
\statementsource{quant-ph.9705052}{pdf-e658890bd7ef-p003-b001}
\statementsource{arXiv.2602.20186}{pdf-585ff62e2309-p003-b014}
\uses{sym_orth}
An erasure set $E \subseteq \mathrm{Fin}\,n$ is \emph{correctable} for $S$ if
every $v \in V_E \cap S^{\perp_\omega}$ is already in $S$.
\end{definition}

\begin{lemma}[Erasure correction below the code distance]
\lean{dist_implies_correctable}
\label{dist_implies_correctable}
\leanok
\proofsource{quant-ph.9705052}{pdf-e658890bd7ef-p003-b001}
\proofsource{arXiv.2602.20186}{
  pdf-585ff62e2309-p003-b015,
  pdf-585ff62e2309-p003-b016
}
\uses{correctable, code_dist}
\difficulty{4}
Let $p$ be prime and let $S$ be an $\mathbb{F}_p$-subspace of the symplectic space $V =
\mathbb{F}_p^n \times \mathbb{F}_p^n$, with code distance $d(S)$ and symplectic
orthogonal complement $S^{\perp_\omega}$. If $E \subseteq \{0,\dots,n-1\}$ is a set of
erased coordinates with $|E| < d(S)$, then $E$ is correctable for $S$: every $v \in V_E
\cap S^{\perp_\omega}$ already lies in $S$.
\end{lemma}

\begin{proof}
\leanok
Any $v \in S^{\perp_\omega} \cap V_E$ outside $S$ witnesses
$d(S) \le \mathrm{wt}(v)$, while $v \in V_E$ forces
$\mathrm{supp}(v) \subseteq E$, so $\mathrm{wt}(v) \le |E| < d(S)$ --- a
contradiction.  Used in the Singleton bound to turn the distance hypothesis
into two disjoint correctable erasure sets of size $d(S)-1$.  Not proved in
Knill--Laflamme (Phys.\ Rev.\ A \textbf{55}, 900); full note:
\texttt{references/informalized/QuantumSingleton.dist\_implies\_correctable.md}.
\end{proof}

\begin{definition}[Logical dimension]
\lean{code_k}
\label{code_k}
\leanok
\statementsource{quant-ph.9705052}{pdf-e658890bd7ef-p004-b008}
\statementsource{arXiv.2602.20186}{pdf-585ff62e2309-p003-b010}
\statementsource{quant-ph.9608006}{
  pdf-4fa1f624d1c5-p002-b006,
  pdf-4fa1f624d1c5-p002-b007
}
\statementsource{quant-ph.0508070}{pdf-4bcdc0ee9c86-p005-b005}
\statementsource{quant-ph.0005008}{
  pdf-1ac7bb3b99cc-p004-b005,
  pdf-1ac7bb3b99cc-p005-b002
}
\statementsource{PhysRevA.55.900}{pdf-bbdd8e5c3949-p004-b002}
\uses{IsIsotropic}
$k(S) = n - \dim(S)$.
\end{definition}

%%%%%%%%%%%%%%%%%%%%%%%%%%%%%%%%%%%%%%%%%%%%%%%%%%%%%%%%%%%%%%%%%%%%%%%%%%%%%%%
\section{Key lemma: two disjoint correctable sets}

\begin{lemma}[Two disjoint correctable erasures bound the logical dimension]
\lean{two_disjoint_correctable_sets_bound_logical_dimension}
\label{two_disjoint_correctable_sets_bound_logical_dimension}
\leanok
\statementsource{quant-ph.9705052}{pdf-e658890bd7ef-p011-b001}
\proofsource{arXiv.2602.20186}{
  pdf-585ff62e2309-p004-b004,
  pdf-585ff62e2309-p004-b005
}
\statementsource{PhysRevA.55.900}{pdf-bbdd8e5c3949-p010-b001}
\uses{correctable, code_k, dim_V_sub, finrank_sym_orth}
\difficulty{4}
Work in the symplectic space $V = \bbf_p^{\,n} \times \bbf_p^{\,n}$ over the prime field
$\bbf_p$, equipped with the symplectic form $\omega$, and let $S \le V$ be an isotropic
submodule, so that $\omega(u,v) = 0$ for all $u, v \in S$. Suppose $A, B \subseteq \{0,
1, \dots, n-1\}$ are two disjoint sets of coordinates, each correctable for $S$: every
vector of $S^{\perp_\omega}$ supported on $A$ lies in $S$, and likewise every vector of
$S^{\perp_\omega}$ supported on $B$ lies in $S$. Then the logical dimension $k(S) = n -
\dim S$ satisfies
\[
k(S) \;\le\; \bigl|\{0, 1, \dots, n-1\} \setminus (A \cup B)\bigr| \;=\; n - |A| - |B|.
\]
\end{lemma}

\begin{proof}
\leanok
By the cleaning identity, $g(S, \mathrm{univ} \setminus A) = 2(n - \dim S)$
(\texttt{g\_complement\_correctable}), while
$g(S, B \cup (\mathrm{univ}\setminus(A \cup B))) \le
2\,|\mathrm{univ}\setminus(A \cup B)|$ (\texttt{g\_le\_two\_card\_C}); the two
argument sets coincide, giving the bound.  This is the engine behind
Theorem~V.1 of Knill--Laflamme (Phys.\ Rev.\ A \textbf{55}, 900), which
states the corresponding bound $n \ge 4e + k$ with proof deferred; the Lean
route is a symplectic dimension count with no density matrices.
\end{proof}

%%%%%%%%%%%%%%%%%%%%%%%%%%%%%%%%%%%%%%%%%%%%%%%%%%%%%%%%%%%%%%%%%%%%%%%%%%%%%%%
\section{Quantum Singleton bound}

\begin{theorem}[Quantum Singleton bound]
\lean{quantum_singleton_bound}
\label{quantum_singleton_bound}
\leanok
\statementsource{quant-ph.9705052}{pdf-e658890bd7ef-p011-b001}
\proofsource{arXiv.2602.20186}{
  pdf-585ff62e2309-p004-b007,
  pdf-585ff62e2309-p004-b008
}
\statementsource{quant-ph.9608006}{
  pdf-4fa1f624d1c5-p012-b005,
  pdf-4fa1f624d1c5-p011-b008
}
\statementsource{quant-ph.0508070}{
  pdf-4bcdc0ee9c86-p008-b009,
  pdf-4bcdc0ee9c86-p008-b010,
  pdf-4bcdc0ee9c86-p009-b001
}
\statementsource{PhysRevA.55.900}{pdf-bbdd8e5c3949-p010-b001}
\uses{two_disjoint_correctable_sets_bound_logical_dimension, dist_implies_correctable, code_dist_le_n}
\difficulty{6}
Let $p$ be prime and work in the symplectic space $V = \bbf_p^{\,n} \times \bbf_p^{\,n}$
equipped with the form $\omega$. If $S \le V$ is an isotropic subspace, so that
$\omega(u,v) = 0$ for all $u, v \in S$, then its logical dimension $k(S) = n - \dim(S)$
and its distance $d(S)$ satisfy
\[
k(S) + 2\bigl(d(S) - 1\bigr) \;\le\; n .
\]
\end{theorem}

\begin{proof}
\leanok
Choose disjoint erasure sets $A, B$ of size $d(S)-1$
(\texttt{exists\_disjoint\_finsets\_card}); both are correctable by
\texttt{dist\_implies\_correctable}, so the two-disjoint-correctable-sets
lemma yields $k(S) \le n - 2(d(S)-1)$.  The statement is Theorem~V.1 of
Knill--Laflamme (Phys.\ Rev.\ A \textbf{55}, 900): an $(n,k)$
$e$-error-correcting code satisfies $n \ge 4e + k$, which with $d = 2e+1$ is
exactly this bound; the paper defers the proof, and the Lean supplies it via
the erasure/cleaning argument.
\end{proof}

%%% added by the blueprint pipeline
\section{Additional declarations}

\begin{definition}[Prime field $\mathbb{F}_p$]
\lean{F}
\label{F}
\leanok
\statementsource{arXiv.2602.20186}{pdf-585ff62e2309-p002-b006}
The prime field $\mathrm{GF}(p)$, realized as $\mathbb{Z}/p\mathbb{Z}$ via \texttt{ZMod}.
\end{definition}

\begin{definition}[Symplectic vector space]
\lean{V}
\statementsource{quant-ph.0005008}{
  pdf-1ac7bb3b99cc-p003-b003,
  pdf-1ac7bb3b99cc-p004-b002
}
\label{V}
\leanok
\statementsource{quant-ph.9705052}{pdf-e658890bd7ef-p007-b003}
\statementsource{arXiv.2602.20186}{pdf-585ff62e2309-p002-b006}
\statementsource{quant-ph.9608006}{
  pdf-4fa1f624d1c5-p002-b005,
  pdf-4fa1f624d1c5-p003-b002
}
\uses{F}
The ambient space $V = \mathbb{F}_p^n \times \mathbb{F}_p^n$ of pairs of coordinate vectors.
\end{definition}

\begin{lemma}[Evaluation of the bundled symplectic form]
\lean{symB_apply}
\label{symB_apply}
\leanok
\uses{V, symB, sym_form}
\difficulty{1}
Fix a prime $p$ and an integer $n$, and set $V = \bbf_p^n \times \bbf_p^n$, equipped
with the alternating bilinear form
\[
\omega\big((x,z),(x',z')\big) \;=\; \sum_{i=1}^{n}\big(x_i z'_i - z_i x'_i\big),
\qquad (x,z),(x',z') \in V.
\]
Let $B \colon V \times V \to \bbf_p$ be the bilinear form on $V$ obtained by bundling
$\omega$. Then $B$ agrees with $\omega$ on every pair of arguments: for all $u, v \in V$
one has $B(u,v) = \omega(u,v)$.
\end{lemma}

\begin{definition}[Restriction to coordinates in $C$]
\lean{restrictToC}
\label{restrictToC}
\leanok
\uses{F, V_sub}
For $C \subseteq \mathrm{Fin}\,n$, the linear map $V_C \to (C \to \mathbb{F}_p) \times (C \to \mathbb{F}_p)$
sending a vector supported on $C$ to its two component functions restricted to $C$.
\end{definition}

\begin{definition}[Extension from coordinates in $C$]
\lean{extendFromC}
\label{extendFromC}
\leanok
\uses{F, V_sub}
The linear map $(C \to \mathbb{F}_p) \times (C \to \mathbb{F}_p) \to V_C$ that extends a pair of
functions on $C$ to a vector of $V$ by setting all coordinates outside $C$ to zero.
\end{definition}

\begin{lemma}[Restriction inverts extension on coordinates in $C$]
\lean{restrictToC_extendFromC}
\label{restrictToC_extendFromC}
\leanok
\uses{extendFromC, restrictToC}
\difficulty{2}
Fix a prime $p$, a natural number $n$, and a subset $C \subseteq \{0,1,\dots,n-1\}$ of
coordinates, and let $V_C$ denote the support submodule of $V = \bbf_p^n \times
\bbf_p^n$ consisting of vectors whose support lies in $C$. Let $E \colon (C \to \bbf_p)
\times (C \to \bbf_p) \to V_C$ be the extension map, which sends a pair $(f,g)$ to the
vector agreeing with $f$ and $g$ on $C$ and equal to zero off $C$, and let $R \colon V_C
\to (C \to \bbf_p) \times (C \to \bbf_p)$ be the restriction map, which sends a vector
supported on $C$ to its two component functions restricted to $C$. Then $R \circ E$ is
the identity: for every pair $x = (f,g)$ of functions $f, g \colon C \to \bbf_p$ one has
$R(E(x)) = x$.
\end{lemma}

\begin{lemma}[Extension recovers a vector supported on $C$]
\lean{extendFromC_restrictToC}
\label{extendFromC_restrictToC}
\leanok
\uses{V_sub, extendFromC, restrictToC}
\difficulty{3}
Fix a prime $p$ and a subset $C \subseteq \{1,\dots,n\}$ of coordinates, and let $V_C
\subseteq \mathbb{F}_p^n \times \mathbb{F}_p^n$ be the submodule of pairs of vectors
whose support is contained in $C$. Then extending back the coordinate functions of a
supported vector recovers that vector exactly: for every $x \in V_C$, applying the
restriction map $V_C \to (C \to \mathbb{F}_p) \times (C \to \mathbb{F}_p)$ and then the
extension-by-zero map back into $V_C$ returns $x$. Equivalently, the extension map is a
left inverse of the restriction map.
\end{lemma}

\begin{definition}[Isomorphism $V_C \cong (\mathbb{F}_p^C)^2$]
\lean{V_sub_iso}
\label{V_sub_iso}
\leanok
\uses{F, V_sub, extendFromC, extendFromC_restrictToC, restrictToC, restrictToC_extendFromC}
The linear equivalence $V_C \simeq (C \to \mathbb{F}_p) \times (C \to \mathbb{F}_p)$ built from
$\mathrm{restrictToC}$ and $\mathrm{extendFromC}$ as mutually inverse maps.
\end{definition}

\begin{definition}[Restriction map $r_E$]
\lean{r_E}
\label{r_E}
\leanok
\statementsource{arXiv.2602.20186}{pdf-585ff62e2309-p003-b009}
\uses{F, V, V_sub}
The linear map $r_E : V \to V_E$ that zeroes out the coordinates of a vector outside $E$,
keeping those in $E$ unchanged.
\end{definition}

\begin{definition}[Intersection $S_M = S \cap V_M$]
\lean{S_M}
\label{S_M}
\leanok
\uses{F, V, V_sub}
For a submodule $S \le V$ and $M \subseteq \mathrm{Fin}\,n$, the submodule $S_M = S \cap V_M$.
\end{definition}

\begin{definition}[Intersection $S^{\perp}_M = S^{\perp_\omega} \cap V_M$]
\lean{S_perp_M}
\label{S_perp_M}
\leanok
\uses{F, V, V_sub, sym_orth}
The submodule $S^{\perp}_M = S^{\perp_\omega} \cap V_M$.
\end{definition}

\begin{definition}[Supportable logical operators count $g(M)$]
\lean{g}
\label{g}
\leanok
\statementsource{arXiv.2602.20186}{pdf-585ff62e2309-p003-b018}
\uses{F, S_M, S_perp_M, V}
$g(S,M) = \dim_{\mathbb{F}_p}(S^{\perp}_M) - \dim_{\mathbb{F}_p}(S_M)$.
\end{definition}

\begin{lemma}[Kernel of the coordinate restriction map]
\lean{ker_r_E}
\label{ker_r_E}
\leanok
\uses{V, V_sub, r_E}
\difficulty{4}
Fix a prime $p$ and let $V = \bbf_p^n \times \bbf_p^n$. For a subset $E \subseteq
\{1,\dots,n\}$, let $r_E \colon V \to V_E$ be the restriction map that keeps the
coordinates indexed by $E$ and sets all others to zero. Then the kernel of $r_E$ is the
support submodule $V_{E^c}$ consisting of those vectors whose support lies in the
complement $E^c = \{1,\dots,n\} \setminus E$; that is, $r_E$ annihilates exactly the
vectors that vanish on every coordinate in $E$.
\end{lemma}

\begin{definition}[Complement of $E$]
\lean{E_c}
\label{E_c}
\leanok
$E^c = \mathrm{Fin}\,n \setminus E$, the complement finset of $E$.
\end{definition}

\begin{lemma}[Complement as set difference]
\lean{E_c_eq}
\label{E_c_eq}
\leanok
\uses{E_c}
\difficulty{1}
Let $E \subseteq \mathrm{Fin}\,n$ be a finite subset. Then its complement equals the set
difference of the whole space and $E$; that is, $E^c = \mathrm{Fin}\,n \setminus E$.
\end{lemma}

\begin{lemma}[Rank of a restricted subspace of $\mathbb{F}_p^n \times \mathbb{F}_p^n$]
\lean{dim_map_r_E}
\label{dim_map_r_E}
\leanok
\uses{E_c, F, V, V_sub, ker_r_E, r_E}
\difficulty{5}
Fix a prime $p$ and an integer $n \ge 0$, and let $V = \mathbb{F}_p^n \times
\mathbb{F}_p^n$. For a set $E$ of coordinates, let $r_E \colon V \to V$ be the
restriction map that leaves each coordinate in $E$ unchanged and sets every coordinate
outside $E$ to zero, and for a set $C$ of coordinates let $V_C \le V$ be the subspace of
vectors supported on $C$. Then for every subspace $S \le V$,
\[
\dim_{\mathbb{F}_p} r_E(S) \;=\; \dim_{\mathbb{F}_p} S \;-\;
\dim_{\mathbb{F}_p}\!\bigl(S \cap V_{E^c}\bigr),
\]
where $E^c$ denotes the complement of $E$ among the $n$ coordinates.
\end{lemma}

\begin{lemma}[Symplectic form is unchanged by restricting the second argument]
\lean{sym_form_r_E}
\label{sym_form_r_E}
\leanok
\uses{V, V_sub, r_E, sym_form}
\difficulty{3}
Fix a prime $p$ and an integer $n$, and work in the space $V = \bbf_p^{\,n} \times
\bbf_p^{\,n}$ equipped with the symplectic form $\omega$. Let $M \subseteq \{0, 1,
\dots, n-1\}$ be a set of coordinates, and let $v \in V_M$, so that every coordinate of
$v$ lying outside $M$ vanishes. Then for every $s \in V$,
\[
\omega(v, s) = \omega(v, r_M(s)),
\]
where $r_M(s)$ is the vector obtained from $s$ by zeroing out all coordinates outside
$M$.
\end{lemma}

\begin{definition}[Restriction as an endomorphism of $V$]
\lean{r_E_V}
\label{r_E_V}
\leanok
\uses{F, V, V_sub, r_E}
The composition $r_E^V : V \to V$ of $r_E$ with the inclusion $V_E \hookrightarrow V$.
\end{definition}

\begin{lemma}[Restriction on the left argument of the symplectic form]
\lean{sym_form_r_E_left}
\label{sym_form_r_E_left}
\leanok
\uses{V, V_sub, r_E_V, sym_form, sym_form_r_E, sym_form_swap}
\difficulty{3}
Let $p$ be prime, $\bbf_p$ the field with $p$ elements, and let $V = \bbf_p^n \times
\bbf_p^n$ carry the symplectic form $\omega\big((x,z),(x',z')\big) =
\sum_{i=0}^{n-1}(x_i z'_i - z_i x'_i)$. For a subset $M \subseteq \{0,\dots,n-1\}$, let
$r_M : V \to V$ be the endomorphism that keeps the coordinates indexed by $M$ and sets
all others to zero, and let $V_M \subseteq V$ be the subspace of vectors whose support
is contained in $M$. Then for every $s \in V$ and every $v \in V_M$, \[
\omega\big(r_M(s),\, v\big) = \omega(s, v). \]
\end{lemma}

\begin{lemma}[Non-degeneracy of the symplectic form on a support submodule]
\lean{sym_form_nondegenerate_on_V_sub}
\label{sym_form_nondegenerate_on_V_sub}
\leanok
\uses{V, V_sub, r_E, sym_form, sym_form_nondegenerate, sym_form_r_E}
\difficulty{3}
Let $p$ be a prime, let $n$ be a natural number, and let $M \subseteq \{0,1,\dots,n-1\}$
be a subset of coordinates. Write $V = \mathbb{F}_p^n \times \mathbb{F}_p^n$ for the
symplectic space equipped with the form $\omega\big((x,z),(x',z')\big) =
\sum_{i=0}^{n-1}(x_i z'_i - z_i x'_i)$, and let $V_M \subseteq V$ be the support
submodule of those vectors whose coordinates outside $M$ all vanish. If $v \in V_M$
satisfies $\omega(v,w) = 0$ for every $w \in V_M$, then $v = 0$.
\end{lemma}

\begin{lemma}[Left restriction is invisible to the symplectic form on the support submodule]
\lean{sym_form_left_restrict}
\label{sym_form_left_restrict}
\leanok
\uses{V, V_sub, r_E, r_E_V, sym_form}
\difficulty{3}
Fix a prime $p$ and a natural number $n$, and let $M \subseteq \{0,1,\dots,n-1\}$ be a
set of coordinates. For every $s \in V$ and every $v \in V_M$ (that is, every $v$
supported on $M$),
\[
\omega\bigl(r_M^V(s),\, v\bigr) = \omega(s, v).
\]
In words: applying the restriction endomorphism $r_M^V$, which zeroes out the
coordinates outside $M$, to the left argument of the symplectic form leaves the value
unchanged whenever the right argument is supported on $M$.
\end{lemma}

\begin{lemma}[Restricting to a support before taking the symplectic complement]
\lean{orth_inter_eq_orth_map}
\label{orth_inter_eq_orth_map}
\leanok
\uses{F, V, V_sub, r_E_V, symB_apply, sym_form, sym_form_left_restrict, sym_orth}
\difficulty{4}
Fix a prime $p$ and work in the symplectic space $V = \bbf_p^{\,n} \times \bbf_p^{\,n}$
equipped with the form $\omega\big((x,z),(x',z')\big) = \sum_{i=0}^{n-1}(x_i z'_i - z_i
x'_i)$, and for a subset $M \subseteq \{0,\dots,n-1\}$ let $V_M$ be the submodule of
vectors whose support lies in $M$ and let $r_M \colon V \to V$ be the map that zeroes
every coordinate outside $M$ and fixes those in $M$. Then for every $\bbf_p$-subspace $S
\le V$,
\[
S^{\perp_\omega} \cap V_M \;=\; \big(r_M(S)\big)^{\perp_\omega} \cap V_M,
\]
where $r_M(S)$ denotes the image of $S$ under $r_M$ and $(\cdot)^{\perp_\omega}$ the
orthogonal complement with respect to $\omega$.
\end{lemma}

\begin{definition}[Restricted bundled form]
\lean{symB_sub}
\label{symB_sub}
\leanok
\uses{F, V_sub, symB}
The symplectic bilinear form $\mathrm{symB}$ restricted to the subspace $V_M$.
\end{definition}

\begin{definition}[Restricted symplectic form]
\lean{sym_form_sub}
\label{sym_form_sub}
\leanok
\uses{F, V_sub, symB_sub}
An abbreviation for $\mathrm{symB\_sub}$, the symplectic form viewed as a bilinear form on $V_M$.
\end{definition}

\begin{lemma}[Evaluation of the restricted symplectic form]
\lean{sym_form_sub_apply}
\label{sym_form_sub_apply}
\leanok
\uses{V, V_sub, sym_form, sym_form_sub}
\difficulty{1}
Fix a prime $p$ and let $\omega$ denote the symplectic form on $V = \bbf_p^{\,n} \times
\bbf_p^{\,n}$. For a subset $M \subseteq \{0, \dots, n-1\}$ and any two vectors $x, y$
in the support submodule $V_M$, the restricted symplectic form evaluated at $(x, y)$
agrees with $\omega(x, y)$ computed on $x$ and $y$ regarded as elements of the ambient
space $V$.
\end{lemma}

\begin{lemma}[Nondegeneracy of the restricted symplectic form]
\lean{sym_form_sub_nondegenerate}
\label{sym_form_sub_nondegenerate}
\leanok
\uses{V_sub, sym_form, sym_form_nondegenerate_on_V_sub, sym_form_sub}
\difficulty{3}
Let $p$ be a prime, and equip $V = \bbf_p^n \times \bbf_p^n$ with the symplectic form
$\omega\big((x,z),(x',z')\big) = \sum_{i=0}^{n-1}(x_i z'_i - z_i x'_i)$. For any subset
$M \subseteq \{0,\dots,n-1\}$, let $V_M$ be the support submodule consisting of those
$(x,z) \in V$ with $x_i = z_i = 0$ for every $i \notin M$. Then the restriction of
$\omega$ to $V_M$ is nondegenerate: if $u \in V_M$ satisfies $\omega(u,v) = 0$ for all
$v \in V_M$, then $u = 0$.
\end{lemma}

\begin{lemma}[Orthogonal complement of a subspace, intersected with a support submodule]
\lean{orth_inter_eq_orth_sub_image}
\label{orth_inter_eq_orth_sub_image}
\leanok
\uses{F, V, V_sub, orth_inter_eq_orth_map, r_E, r_E_V, symB, sym_form_sub, sym_orth}
\difficulty{3}
Let $p$ be a prime and let $\omega$ be the symplectic form on $V = \bbf_p^{\,n} \times
\bbf_p^{\,n}$. For any subset $M$ of $\{0, 1, \dots, n-1\}$ and any submodule $S \le V$,
the vectors of the support submodule $V_M$ that are $\omega$-orthogonal to every element
of $S$ are exactly the images, under the inclusion $V_M \hookrightarrow V$, of the
vectors of $V_M$ that are orthogonal to $r_M(S)$ with respect to the restriction of
$\omega$ to $V_M$. In symbols, writing $\iota \colon V_M \hookrightarrow V$ for the
inclusion and $\perp_M$ for orthogonality with respect to the restricted form,
\[
S^{\perp_\omega} \cap V_M \;=\; \iota\!\left( r_M(S)^{\perp_M} \right),
\]
where $r_M(S)$ is the image of $S$ under the restriction map $r_M \colon V \to V_M$.
\end{lemma}

\begin{lemma}[Reflexivity of the restricted symplectic form]
\lean{sym_form_sub_isRefl}
\label{sym_form_sub_isRefl}
\leanok
\uses{V, sym_form, sym_form_sub, sym_form_sub_apply, sym_form_swap}
\difficulty{3}
Fix a prime $p$, a natural number $n$, and a subset $M \subseteq \{0,1,\dots,n-1\}$, and
let $V_M$ be the associated support submodule of $V = \bbf_p^{\,n} \times \bbf_p^{\,n}$.
Then the restriction to $V_M$ of the symplectic form $\omega(u,v) = \sum_{i} (x_i z'_i -
z_i x'_i)$, where $u=(x,z)$ and $v=(x',z')$, is reflexive: for all $u, v \in V_M$, if
$\omega(u,v) = 0$ then $\omega(v,u) = 0$.
\end{lemma}

\begin{lemma}[Dimension of the orthogonal intersection over a coordinate set]
\lean{dim_orth_inter}
\label{dim_orth_inter}
\leanok
\uses{F, V, V_sub, dim_V_sub, orth_inter_eq_orth_sub_image, r_E, sym_form_sub, sym_form_sub_isRefl, sym_form_sub_nondegenerate, sym_orth}
\difficulty{4}
Fix a prime $p$, an integer $n$, and work in the symplectic space $V = \bbf_p^{\,n}
\times \bbf_p^{\,n}$ equipped with the symplectic form $\omega\big((x,z),(x',z')\big) =
\sum_{i} (x_i z'_i - z_i x'_i)$. Let $M$ be a subset of the coordinate set
$\{0,1,\dots,n-1\}$, let $V_M \le V$ be the support submodule of vectors whose
coordinates vanish outside $M$, and let $r_M : V \to V_M$ be the restriction map that
zeroes every coordinate outside $M$. Then for any $\bbf_p$-subspace $S \le V$, the
symplectic orthogonal complement $S^{\perp_\omega}$ satisfies
\[
\dim_{\bbf_p}\big(S^{\perp_\omega} \cap V_M\big) \;=\; 2\,\abs{M} \;-\;
\dim_{\bbf_p}\big(r_M(S)\big).
\]
\end{lemma}

\begin{lemma}[Expansion of the operator count $g(S,M)$]
\lean{g_expansion}
\label{g_expansion}
\leanok
\uses{E_c, F, IsIsotropic, S_M, S_perp_M, V, V_sub, dim_V_sub, dim_map_r_E, dim_orth_inter, g, r_E, sym_orth}
\difficulty{5}
Fix a prime $p$ and equip $V = \bbf_p^{\,n} \times \bbf_p^{\,n}$ with the symplectic
form $\omega\big((x,z),(x',z')\big) = \sum_{i=0}^{n-1}\big(x_i z'_i - z_i x'_i\big)$.
For a set of coordinates $M \subseteq \{0,\dots,n-1\}$, let $V_M \le V$ be the subspace
of vectors supported on $M$, and for a subspace $W \le V$ write $W_M = W \cap V_M$. Let
$S \le V$ be isotropic, meaning $\omega(u,v) = 0$ for all $u,v \in S$, let $S^{\perp}$
denote its symplectic orthogonal complement, and set $g(S,M) =
\dim_{\bbf_p}\!\big(S^{\perp}\cap V_M\big) - \dim_{\bbf_p}(S_M)$. Then, writing $M^c$
for the complementary set of coordinates,
\[
g(S,M) \;=\; 2\,\abs{M} + \dim_{\bbf_p}\!\big(S_{M^c}\big) - \dim_{\bbf_p}(S) -
\dim_{\bbf_p}(S_M).
\]
\end{lemma}

\begin{lemma}[Double complement in a finite set]
\lean{E_c_E_c}
\label{E_c_E_c}
\leanok
\uses{E_c}
\difficulty{1}
Fix a natural number $n$, and for a subset $E \subseteq \mathrm{Fin}\,n$ write $E^c =
\mathrm{Fin}\,n \setminus E$ for its complement. Then for every subset $M \subseteq
\mathrm{Fin}\,n$, taking the complement twice recovers the original set: $(M^c)^c = M$.
\end{lemma}

\begin{lemma}[Additivity bound for support-restricted subspaces]
\lean{dim_S_M_add_dim_S_M_c_le_dim_S}
\label{dim_S_M_add_dim_S_M_c_le_dim_S}
\leanok
\uses{E_c, F, S_M, V, V_sub}
\difficulty{4}
Fix a prime $p$ and an integer $n$, and let $S$ be a subspace of the ambient space $V =
\mathbb{F}_p^n \times \mathbb{F}_p^n$. For a subset $M \subseteq \{1,\dots,n\}$ with
complement $M^c$, write $S_M = S \cap V_M$ and $S_{M^c} = S \cap V_{M^c}$ for the
intersections of $S$ with the corresponding support subspaces. Then
\[
  \dim_{\mathbb{F}_p} S_M + \dim_{\mathbb{F}_p} S_{M^c} \le \dim_{\mathbb{F}_p} S.
\]
\end{lemma}

\begin{lemma}[Dimension formula for the operator count of an isotropic subspace]
\lean{g_formula}
\label{g_formula}
\leanok
\uses{E_c, F, IsIsotropic, S_M, V, g, g_expansion}
\difficulty{2}
Let $p$ be prime, and equip the symplectic space $V = \mathbb{F}_p^n \times
\mathbb{F}_p^n$ with its standard alternating form $\omega$. If $S \le V$ is an
isotropic subspace and $M \subseteq \{0,1,\dots,n-1\}$ is any set of coordinates, then
\[
g(S,M) \;=\; \bigl(2\,\abs{M} + \dim_{\mathbb{F}_p} S_{M^c}\bigr) \;-\;
\bigl(\dim_{\mathbb{F}_p} S + \dim_{\mathbb{F}_p} S_M\bigr),
\]
where $M^c$ is the complement of $M$ and $S_{M^c} = S \cap V_{M^c}$ consists of the
vectors of $S$ supported on $M^c$.
\end{lemma}

\begin{lemma}[Additive dimension identity for an isotropic subspace]
\lean{g_add_dims}
\label{g_add_dims}
\leanok
\uses{E_c, F, IsIsotropic, S_M, V, V_sub, dim_S_M_add_dim_S_M_c_le_dim_S, dim_V_sub, dim_map_r_E, dim_orth_inter, g, g_formula, r_E, sym_orth}
\difficulty{5}
Let $S$ be an isotropic submodule of the symplectic space $V=\bbf_p^n\times\bbf_p^n$,
and let $M\subseteq\{0,\dots,n-1\}$ be a set of coordinates with complement $M^c$.
Writing $S_M=S\cap V_M$ and $S_{M^c}=S\cap V_{M^c}$ for the parts of $S$ supported on
$M$ and on its complement, one has
\[
g(S,M)+\dim_{\bbf_p}(S_M)+\dim_{\bbf_p}(S) \;=\; 2\lvert M\rvert +
\dim_{\bbf_p}(S_{M^c}).
\]
\end{lemma}

\begin{lemma}[A dimension inequality for isotropic subspaces]
\lean{dim_ineq_aux}
\label{dim_ineq_aux}
\leanok
\uses{E_c, F, IsIsotropic, S_M, V, g_add_dims}
\difficulty{1}
Let $p$ be a prime and $V = \bbf_p^n \times \bbf_p^n$, equipped with the symplectic form
$\omega\big((x,z),(x',z')\big) = \sum_{i} (x_i z'_i - z_i x'_i)$. For a coordinate
subset $M \subseteq \{0,\dots,n-1\}$, let $V_M \le V$ be the support submodule of
vectors supported on $M$, and for a subspace $S \le V$ write $S_M = S \cap V_M$. If $S$
is isotropic, that is $\omega(u,v) = 0$ for all $u,v \in S$, then \[ \dim_{\bbf_p} S +
\dim_{\bbf_p} S_M \;\le\; 2\,\abs{M} + \dim_{\bbf_p} S_{M^c}, \] where $M^c =
\{0,\dots,n-1\} \setminus M$ is the complementary coordinate set.
\end{lemma}

\begin{lemma}[Cleaning dimension identity for isotropic subspaces]
\lean{cleaning_dimension_identity}
\label{cleaning_dimension_identity}
\leanok
\proofsource{arXiv.2602.20186}{
  pdf-585ff62e2309-p003-b019,
  pdf-585ff62e2309-p003-b020
}
\uses{E_c, E_c_E_c, F, IsIsotropic, S_M, V, g, g_add_dims}
\difficulty{4}
Let $p$ be prime and let $S \le V = \bbf_p^n \times \bbf_p^n$ be a subspace that is
isotropic for the symplectic form $\omega$, meaning $\omega(u,v)=0$ for all $u,v \in S$.
Then for every subset $M \subseteq \{0,\dots,n-1\}$, with complement $M^c$,
\[
g(S,M) + g(S,M^c) = 2n - 2\dim_{\bbf_p} S,
\]
where $g(S,M) = \dim_{\bbf_p}\!\big(S^{\perp}_M\big) - \dim_{\bbf_p}\!\big(S_M\big)$.
\end{lemma}

\begin{lemma}[Cardinality of a subset plus its complement]
\lean{card_add_compl}
\label{card_add_compl}
\leanok
\uses{E_c}
\difficulty{2}
Let $n$ be a natural number and let $M$ be a subset of the $n$-element index set
$\mathrm{Fin}\,n = \{0, 1, \dots, n-1\}$. Then the size of $M$ and the size of its
complement $M^c = \mathrm{Fin}\,n \setminus M$ add up to $n$: \[ \abs{M} + \abs{M^c} =
n. \]
\end{lemma}

\begin{lemma}[Correctable coordinate sets carry no supportable logical operators]
\lean{correctable_implies_g_zero}
\label{correctable_implies_g_zero}
\leanok
\proofsource{arXiv.2602.20186}{pdf-585ff62e2309-p004-b002}
\uses{F, S_M, V, V_sub, correctable, g, sym_orth}
\difficulty{2}
Fix a prime $p$, let $V = \bbf_p^n \times \bbf_p^n$ be equipped with the symplectic form
$\omega\big((x,z),(x',z')\big) = \sum_{i} (x_i z'_i - z_i x'_i)$, and let $S \le V$ be a
subspace with symplectic orthogonal complement $S^{\perp_\omega} = \{v \in V :
\omega(v,s) = 0 \text{ for all } s \in S\}$. Let $M \subseteq \{0,1,\dots,n-1\}$ be a
set of coordinates, write $V_M$ for the subspace of vectors supported on $M$, and set
$S_M = S \cap V_M$ and $S^{\perp}_M = S^{\perp_\omega} \cap V_M$. If $M$ is correctable
for $S$ — meaning every vector of $S^{\perp_\omega}$ supported on $M$ already lies in
$S$, i.e. $S^{\perp}_M \le S$ — then $g(S,M) = \dim_{\bbf_p}(S^{\perp}_M) -
\dim_{\bbf_p}(S_M) = 0$.
\end{lemma}

\begin{lemma}[Correctable erasure and the count of supportable logical operators]
\lean{g_complement_correctable}
\label{g_complement_correctable}
\leanok
\proofsource{arXiv.2602.20186}{pdf-585ff62e2309-p004-b002}
\uses{E_c, F, IsIsotropic, V, cleaning_dimension_identity, correctable, correctable_implies_g_zero, g}
\difficulty{2}
Let $S$ be an isotropic subspace of the symplectic space $V = \mathbb{F}_p^n \times
\mathbb{F}_p^n$ (with $p$ prime), and let $M \subseteq \{0,1,\dots,n-1\}$ be an erasure
set that is correctable for $S$. Writing $M^c$ for the complement of $M$, one has
\[
g(S, M^c) \;=\; 2n - 2\dim_{\mathbb{F}_p} S.
\]
\end{lemma}

\begin{lemma}[Correctable part contributes nothing to the operator count]
\lean{g_le_two_card_C}
\label{g_le_two_card_C}
\leanok
\proofsource{arXiv.2602.20186}{pdf-585ff62e2309-p004-b005}
\uses{F, S_M, S_perp_M, V, V_sub, correctable, dim_V_sub, g, r_E, sym_orth}
\difficulty{6}
Let $p$ be prime and let $S$ be an $\mathbb{F}_p$-subspace of the symplectic space $V =
\mathbb{F}_p^n \times \mathbb{F}_p^n$. Suppose $B$ and $C$ are disjoint subsets of the
coordinate set $\{0,\dots,n-1\}$, and that $B$ is correctable for $S$. Then
\[
g(S,\, B \cup C) \;\le\; 2\lvert C\rvert,
\]
where $g(S,M) = \dim_{\mathbb{F}_p}\!\bigl(S^{\perp}_M\bigr) -
\dim_{\mathbb{F}_p}\!\bigl(S_M\bigr)$.
\end{lemma}

\begin{proof}
\leanok
Rank--nullity for the restriction map
$r_C : S^{\perp_\omega} \cap V_{B \cup C} \to V_C$: its kernel consists of
vectors supported on $B$, which correctability of $B$ absorbs into $S$, and
its range has dimension at most $\dim V_C = 2|C|$.  This is the counting half
of the cleaning argument, consumed by the two-disjoint-correctable-sets
lemma.  Not present in Knill--Laflamme (Phys.\ Rev.\ A \textbf{55}, 900);
full note:
\texttt{references/informalized/QuantumSingleton.g\_le\_two\_card\_C.md}.
\end{proof}

\begin{lemma}[Full index set gives the whole space]
\lean{V_sub_univ_eq_top}
\label{V_sub_univ_eq_top}
\leanok
\uses{F, V, V_sub}
\difficulty{2}
Fix a prime $p$ and an integer $n$, and let $V = \bbf_p^n \times \bbf_p^n$. Taking $C$
to be the full index set $\{0,1,\dots,n-1\}$, the support submodule $V_C = \{v \in V
\mid \mathrm{supp}(v) \subseteq C\}$ is all of $V$; that is, $V_C = V$.
\end{lemma}

\begin{lemma}[Dimension of the ambient space]
\lean{finrank_V}
\label{finrank_V}
\leanok
\uses{F, V, V_sub, V_sub_univ_eq_top, dim_V_sub}
\difficulty{3}
Let $p$ be a prime and $\bbf_p$ the field with $p$ elements, and for a natural number
$n$ set $V = \bbf_p^n \times \bbf_p^n$. Then $V$ is a vector space over $\bbf_p$ of
dimension \[ \dim_{\bbf_p}(V) = 2n. \]
\end{lemma}

\begin{lemma}[Non-degeneracy of the standard symplectic form]
\lean{symB_nondegenerate}
\label{symB_nondegenerate}
\leanok
\proofsource{arXiv.2602.20186}{
  pdf-585ff62e2309-p003-b001,
  pdf-585ff62e2309-p003-b002
}
\uses{symB, symB_apply, sym_form_nondegenerate}
\difficulty{2}
Fix a prime $p$ and an integer $n$, and let $V = \bbf_p^n \times \bbf_p^n$ carry the
bilinear form
\[
\omega\big((x,z),(x',z')\big) \;=\; \sum_{i=1}^{n}\big(x_i z'_i - z_i x'_i\big)
\]
over the prime field $\bbf_p$. Then $\omega$ is non-degenerate: the only vector $u \in
V$ with $\omega(u,v) = 0$ for every $v \in V$ is $u = 0$.
\end{lemma}

\begin{lemma}[Reflexivity of the symplectic form]
\lean{symB_isRefl}
\label{symB_isRefl}
\leanok
\uses{symB, symB_apply, sym_form, sym_form_swap}
\difficulty{2}
Let $p$ be a prime, $n$ a natural number, and let $\omega$ be the bilinear form on $V =
\bbf_p^n \times \bbf_p^n$ given by $\omega(u,v) = \sum_{i=1}^{n}(x_i z'_i - z_i x'_i)$
for $u = (x,z)$ and $v = (x',z')$. Then $\omega$ is reflexive: for all $u, v \in V$, if
$\omega(u,v) = 0$ then $\omega(v,u) = 0$.
\end{lemma}

\begin{lemma}[Dimension bound for isotropic subspaces]
\lean{finrank_le_n_of_isotropic}
\label{finrank_le_n_of_isotropic}
\leanok
\statementsource{quant-ph.9608006}{pdf-4fa1f624d1c5-p003-b004}
\uses{F, IsIsotropic, V, finrank_V, finrank_sym_orth, sym_orth}
\difficulty{4}
Let $p$ be a prime and equip the $2n$-dimensional space $V = \bbf_p^n \times \bbf_p^n$
over $\bbf_p$ with the symplectic form $\omega\big((x,z),(x',z')\big) = \sum_{i=1}^{n}
(x_i z'_i - z_i x'_i)$. If $S \le V$ is an isotropic subspace, meaning $\omega(u,v) = 0$
for all $u, v \in S$, then $\dim_{\bbf_p} S \le n$.
\end{lemma}

\begin{lemma}[Zero logical dimension forces zero distance]
\lean{code_dist_eq_zero_of_code_k_eq_zero}
\label{code_dist_eq_zero_of_code_k_eq_zero}
\leanok
\uses{F, IsIsotropic, V, code_dist, code_k, finrank_le_n_of_isotropic, finrank_sym_orth, sym_orth, wt}
\difficulty{5}
Let $p$ be a prime and let $S$ be a subspace of the symplectic space $V = \bbf_p^n
\times \bbf_p^n$. If $S$ is isotropic and its logical dimension vanishes, $k(S) = 0$,
then the code distance vanishes as well, $d(S) = 0$.
\end{lemma}

\begin{lemma}[Existence of two disjoint sets of equal size]
\lean{exists_disjoint_finsets_card}
\label{exists_disjoint_finsets_card}
\leanok
\difficulty{4}
Let $n$ and $t$ be natural numbers with $2t \le n$. Then there exist two disjoint
subsets $A, B \subseteq \{1, \dots, n\}$, each of cardinality $t$.
\end{lemma}
